\subsection{Confirmatory PJM decision outcome}
The robust gate admitted a learned model in all four regions. The frozen decisions selected ridge regression for AEP, COMED, and PJME, and a radius-controlled one-rule TSK model for DAYTON. Relative to the seven-day seasonal-naive fallback, independent test RMSE improved by 57.97\%, 55.40\%, 49.13\%, and 51.70\%, respectively; the corresponding asymmetric-cost improvements were 58.64\%, 54.11\%, 50.30\%, and 51.42\%. Mean improvements were 53.55\% in RMSE and 53.62\% in cost. Selected certificate values ranged from 0.0841 to 0.0959, and test target and forecast clipping were zero.

This result supports the transfer of the redesigned deployment gate after the negative Tetouan development case, but it is not evidence that a nonlinear fuzzy model was responsible for the improvement. Three decisions were ridge models, and the selected DAYTON TSK contained one rule and was therefore an affine predictor. No fixed-eight-rule TSK candidate passed the frozen gate. The confirmatory case consequently supports complexity-aware abstention and selection, while the specific structural-sparsity claim remains grounded in the separate fixed-$K$ versus radius-controlled ablation. The comparison is also limited to the predeclared seasonal-naive operational fallback and does not establish dominance over all forecasting methods.
